\section{Implementation}\label{sec:implementation}

The method is implemented as the Python package \code{motiongate} (Python $\ge3.10$, NumPy and OpenCV; Ultralytics only for the detector wrapper). Its modules mirror \cref{sec:method}: \code{gates} provides \FD, \FDS{} and MOG2 behind one interface, \code{scheduler} implements rule~\eqref{eq:rule}, \code{detector} wraps Ultralytics (any callable that returns boxes can replace it), \code{theory} contains the closed-form predictions of \cref{sec:envelope}, and \code{cli} is the command-line tool. A minimal example:

\begin{lstlisting}
from motiongate import (MotionGatedDetector, SchedulerConfig,
                        UltralyticsPersonDetector, make_gate)

system = MotionGatedDetector(
    UltralyticsPersonDetector("yolov8m.pt"),
    make_gate("stabilized"),
    SchedulerConfig.from_max_staleness(seconds=0.5, fps=30))
for frame in frames:                        # BGR uint8 images
    result = system.process(frame)
    print(result.index, result.reason, result.detection_age, result.detections)
\end{lstlisting}

\noindent The command-line tool processes a video file or a camera and writes an annotated video and a per-frame CSV log, e.g.\ \code{motiongate input.mp4 --max-staleness 0.5 --output out.mp4 --log frames.csv}.

Several design decisions matter in deployment. Input frames are never modified. A change of resolution restarts the stream, because the gate compares aligned frames. Every result states whether its boxes are fresh or retained, why the detector did or did not run (initial frame, motion, refresh, skipped) and how old the boxes are, so that downstream logic can, for instance, treat stale boxes differently. The package comes with 28 unit tests covering the behaviour of the gates on synthetic textures (static scenes, moving objects, translations, illumination steps, noise), the guarantees of Proposition~\ref{prop:guarantees}, decision-by-decision equivalence of the online scheduler and the offline replica used in the experiments, and agreement of the noise model~\eqref{eq:bias} with the actual OpenCV pipeline.
